% =====================================================================
% VI. LIMITATIONS AND CONCLUSION -- OWNER: Conor
% =====================================================================
\section{Limitations and Conclusion}\label{sec:conclusion}
% OWNER: Conor. STARTER DRAFT (written with Claude Code, 2026-10-20): Conor to
% check every sentence against its source, rewrite, and commit as his own.
% Sources: outputs/clean/paper_stats.json; journal/2026-10-07-circuit-tracing.md;
% journal/2026-10-16.md; 50-problem report "Data integrity"
\textbf{Limitations.} We test one model on one dataset, with 50 problems per mode, so
the intervals are wide. The checker sees only calculations written as equations or
with operator words. A step such as ``she has 9 left'' is invisible to it, and 17 of
the 100 responses could not be compared at all. Its coverage, estimated from hand
counts, is \textbf{TBD}. Steps within a response are not independent, so the
step-level interval is optimistic. The planted-working and copy-vs-use tests use
Qwen3-0.6B, while the GSM8K outputs come from 1.7B. Circuit-tracing graphs describe a
transcoder replacement of the model, not Qwen3 itself. Finally, the FOLIO results
validate the logic checker only.

\textbf{Future work.} Natural errors are too rare on GSM8K to measure unfaithfulness
directly. Harder problems, where working and answers diverge more often, would give
the test more power. Translating free-text reasoning into first-order logic, with
human review, would let us check Qwen3 on FOLIO. Feature interventions on attribution
graphs could test whether the answer depends on the working.
% Add the hint test / copy-vs-use here if either missed the Wed 21 Oct 9pm freeze.

\textbf{Conclusion.} On GSM8K, Qwen3's final answers never contradicted its own
written arithmetic, with or without thinking. That shows the model is consistent, not
that it is faithful. Telling the two apart needs tests that change the reasoning or
the prompt, and our results from those are \textbf{TBD}.
